\documentclass[10pt,a4paper]{amsart}
\usepackage[margin=19mm]{geometry}
\usepackage{amsmath,amssymb,mathtools}
\usepackage[scr=rsfs,cal=euler]{mathalfa}
\usepackage{xfrac}
\usepackage[unicode,hidelinks,bookmarks=false]{hyperref}
\newcommand{\ie}{i.e.\ }
\newcommand{\cf}{cf.\ }
\newcommand{\resp}{respectively\ }
\newcommand{\Subsection}{\subsection}
\providecommand{\nobreakdash}{\nobreak\hbox{-}}
\setlength{\emergencystretch}{2em}
\multlinegap0pt
\usepackage{bm}
\usepackage{bbm}
\usepackage{dsfont}
\usepackage{xspace}
\usepackage{cancel}
\usepackage{mathtools}
\usepackage{amsthm}
\makeatletter
\@ifundefined{thm}{\newtheorem{thm}{Thm}}{}
\@ifundefined{prop}{\newtheorem{prop}[thm]{Proposition}}{}
\makeatother
\newcommand{\M}{\mathcal{M}}
\newcommand{\R}{\mathbb{R}}
\title[Smooth globally PLI functions are nonlinear least-squares, a]{Smooth globally PLI functions are nonlinear least-squares, and so are their gradient-dominated cousins}
\author{Eduardo D. Sontag}
\date{Source: arXiv:2608.08849v2 (CC BY 4.0)}
\begin{document}
\maketitle

\noindent\emph{Source and licence.} Smooth globally PLI functions are nonlinear least-squares, and so are their gradient-dominated cousins. Version arXiv:2608.08849v2 (CC BY 4.0), distributed under the Creative Commons Attribution 4.0 International licence, as recorded in the arXiv licence metadata for this exact version and shown on the abstract page https://arxiv.org/abs/2608.08849v2. Full rights notice: licenses/arxiv-2608.08849-CC-BY-4.0.txt.\par\smallskip

\noindent\textit{Editorial scope.} Counted unit: the Prop whose source label is \texttt{prop:sgl} in the article (source0.tex lines 360--367, source label \texttt{prop:sgl}), delivered together with the article's own complete original proof of it (source0.tex lines 369--405, one \texttt{proof} environment of 37 lines). No mathematics is written, repaired or re-derived: statement and proof are the article's own text line for line. Required context retained uncounted: eq:GD (lines 257--260); eq:A (lines 266--269). Every definition, display and label of those lines is retained unchanged. Omitted from this package and not redistributed: 1--256, 261--265, 270--359, 368--368, 406--2562. Nothing inside a retained excerpt is omitted or altered: every retained body is delivered exactly as the article's lines are written, so no omission receipt of any kind accompanies this package. Citations: the retained excerpts carry no citation command at all, so no bibliography entry of the article is transcribed and no reference is redistributed. Reference adaptation: none was needed and none was made; every reference inside the delivered unit resolves inside it, so no unresolved reference or citation is produced. Printed slips kept literal: no printed word, symbol, hypothesis or number of the counted statement or of its proof was altered. Layout and class: the article's own class is \texttt{article} with packages including amsmath, amssymb, amsthm, newtxtext, newtxmath, fontenc, inputenc, geometry, xcolor, graphicx, microtype, hyperref; the wrapper is the lane's pinned \texttt{amsart} editorial wrapper at 10pt on A4, which restates the theorem-like environments the retained text needs and adds the article's own macro definitions that the retained text needs (\texttt{\textbackslash M}, \texttt{\textbackslash R}), with extra wrapper packages (bm, bbm, dsfont, xspace, cancel, mathtools). The delivered numbering is the unit's own. Assets: no retained span contains a figure environment, a graphics-inclusion command, a PDF-page inclusion, a table or a TikZ picture; no image file is copied into this package, the package declares no asset dependency and no image file is redistributed with it. Licence: version arXiv:2608.08849v2 of this article is distributed under the Creative Commons Attribution 4.0 International licence, as recorded in the arXiv licence metadata for this exact version and shown on the abstract page https://arxiv.org/abs/2608.08849v2. A full rights notice is delivered at licenses/arxiv-2608.08849-CC-BY-4.0.txt. Characters: every retained excerpt is pure ASCII, so the delivered engine, XeLaTeX with its default Latin Modern text font, reports no missing character.
\begin{equation}
  \|\nabla f(x)\| \ \ge\ \alpha\bigl(h(x)\bigr) \qquad \text{for all } x \in \M .
  \tag{$\mathrm{GD}_\alpha$}\label{eq:GD}
\end{equation}

\begin{equation}
  \alpha(s)^2 \ \ge\ 2\mu s \qquad \text{for all } s \in [0,h_0] .
  \tag{A}\label{eq:A}
\end{equation}

\begin{prop}[Our hypothesis is $sgl$-P\L{}I]\label{prop:sgl}
Let $f\colon\M\to\R$ be $C^1$ and bounded below. The following are equivalent:
\begin{enumerate}
\item[(i)] $f$ satisfies \eqref{eq:GD} for some positive definite $\alpha$
satisfying \eqref{eq:A};
\item[(ii)] $f$ satisfies $sgl$-P\L{}I.
\end{enumerate}
\end{prop}

\begin{proof}
(i)$\Rightarrow$(ii). Fix $\rho>0$; we may assume $\rho>h_0$. On $\{h\le h_0\}$,
\eqref{eq:A} gives $\|\nabla f\|^2\ge\alpha(h)^2\ge 2\mu h$. On
$\{h_0\le h\le\rho\}$, let $m:=\min_{[h_0,\rho]}\alpha>0$, which is positive
because $\alpha$ is continuous and strictly positive on the compact interval
$[h_0,\rho]$; then $\|\nabla f\|^2\ge m^2\ge (m^2/\rho)h$. So
$c_\rho:=\min\{2\mu, m^2/\rho\}$ works.

(ii)$\Rightarrow$(i). Here we must exhibit a single $\alpha$ that is positive
definite, satisfies \eqref{eq:GD}, and satisfies \eqref{eq:A}.

Dispose first of the constant case: if $h\equiv0$ then \eqref{eq:GD} holds for
$\alpha(s)=\sqrt s$, which is positive definite and satisfies \eqref{eq:A}. (This
case must be separated out, because for constant $f$ every positive constant is
admissible in (ii) and the supremum below is infinite.) So assume $f$ is not
constant.

Fix any $h_0>0$, and let $c(\rho)$ denote the supremum of the constants
admissible in (ii) for a given $\rho$; a constant admissible for $\rho$ is
admissible for every smaller value, so $c$ is nonincreasing, and it is positive
and finite by hypothesis. Put
\[
  \gamma(s) \ :=\ \int_1^2 c(sv)\,dv \ =\ \frac1s\int_s^{2s} c(u)\,du , \qquad
  \tilde\gamma(s) \ :=\ \min\{\gamma(s), \gamma(h_0)\} .
\]
Then $\gamma$ is nonincreasing (because $c$ is), continuous (the second
expression exhibits it as an average of a monotone, hence locally integrable,
function over an interval varying continuously with $s$), and satisfies
$\gamma\le c$ pointwise. Hence $\tilde\gamma$ is continuous, positive, bounded by
$\gamma(h_0)$, and $\le c$. Setting
$\alpha(s):=\sqrt{\tilde\gamma(s)\,s}$ gives a continuous function with
$\alpha(s)\le\sqrt{\gamma(h_0)s}\to0$ as $s\downarrow0$, so $\alpha(0)=0$; it is
strictly positive for $s>0$, hence positive definite; it satisfies \eqref{eq:GD}
since $\|\nabla f\|^2\ge c(h)h\ge\tilde\gamma(h)h=\alpha(h)^2$; and for
$s\le h_0$ monotonicity of $\gamma$ gives $\tilde\gamma(s)=\gamma(h_0)$, so
\eqref{eq:A} holds with $2\mu=\gamma(h_0)$.
\end{proof}

\end{document}
